\begin{table}[tbp]
\centering
\caption{研究框架、技术实现与解释边界}
\label{tab:model-framework}
\small
\begin{threeparttable}
\begin{tabularx}{\textwidth}{p{0.16\textwidth}p{0.26\textwidth}Y}
\toprule
模块 & 技术实现 & 本文用途与解释边界 \\
\midrule
压力指数构造 & LSI 与 MarketLSI & 将分钟级 OHLCV 与成交额信息压缩为短时压力代理指标；不等同于订单簿真实流动性。 \\
动态风险刻画 & \mbox{RGARCH-CARR-SK} & 刻画 MarketLSI 压力创新的条件风险和高阶矩路径；用于风险度量而非收益率预测。 \\
尾部分位传导 & QVAR & 描述市场状态变量在不同分位点下的联动和情景响应；情景模拟不作严格因果识别。 \\
样本外预警 & SMARTBoost & 检验无泄漏历史特征对未来 H5/H10 压力事件的预测能力；作为唯一正式机器学习预警模型。 \\
\bottomrule
\end{tabularx}
\begin{tablenotes}[flushleft]
\footnotesize
\item 注：三类模型形成递进证据链，而不是并列的模型竞赛。
\end{tablenotes}
\end{threeparttable}
\end{table}
